%% EXHALE (EXHALE): boundary and initial conditions.
%% Complete description of the hydrodynamic boundary conditions
%% (Apply_BC.f90), the "breathing base" phenomenology, the boundary
%% treatment of the ionization/temperature/metal fields, the
%% boundary-adjacent settings (radiative-transfer integration, escape
%% radius, domain truncation, Newton-solver interaction), and the
%% three initial-condition families (set_IC.f90: cold hydrostatic,
%% transonic isothermal wind, hot-Parker warm seed).
%% Updated 2026-08-15: runtime "Grid cells:", the base grid resolution
%% key, the ghost-pressure closures (Hydrostatic base / Base ghost
%% temperature), Low-Mach contact-mode damping, the composition module as
%% the single source of rho_bc / ntot_bc, the residual diagnostic below
%% the escape radius, the restart guards, and the du-trigger arming.
%% Author: Kwang-Il Seon
\documentclass[english,a4paper]{my_memo}
\synctex=1
\usepackage{listings}
\usepackage{xcolor}
\usepackage{booktabs}
\usepackage{amsmath}
\usepackage{pdflscape}
\usepackage{babel}

\newcommand{\Y}{\checkmark}
\newcommand{\N}{--}
\emergencystretch=3em

% listing style for code snippets
\lstdefinestyle{code}{
  basicstyle=\ttfamily\small,
  commentstyle=\color{gray}\itshape,
  frame=single,
  breaklines=true,
  columns=fullflexible,
  showstringspaces=false,
}
\newcommand{\rp}{$R_p$}
\title{Boundary and initial conditions in EXHALE (EXHALE)}
\author{Kwang-Il Seon}
\date{2026-06-12 (updated 2026-08-19)}

\begin{document}
\maketitle

\begin{quote}\small\itshape
This memo documents the boundary conditions and the initial conditions
of the EXHALE / EXHALE hydrodynamics solver at the level of the
actual implementation: what is imposed, where in the code, why it has
the form it does, and what dynamical behavior it produces (in
particular the ``breathing base''). File references are relative to
the repository root; the structure is inherited from ATES v2
(Caldiroli et al.) unless noted as an EXHALE extension.
\end{quote}

\tableofcontents

\section{Overview: where and when the BCs act}
\label{sec:overview}

\subsection{Ghost-cell layout}

The radial grid carries $N$ computational cells ($j = 1,\dots,N$) plus
$N_g = 2$ ghost cells on each side, so all state arrays are dimensioned
\verb|(1-Ng : N+Ng)|, i.e.\ $j = -1, 0$
below the base and $j = N{+}1, N{+}2$ above the top
(\texttt{src/modules/init/parameters.f90}, \texttt{define\_grid.f90}).

\paragraph{2026-08-15: the cell count is a runtime quantity.}
$N$ used to be a compile-time \texttt{parameter} fixed at 500. It is now
the module variable \texttt{integer :: N = 500} of
\texttt{parameters.f90}, set by the optional \texttt{input.inp} key
\texttt{Grid cells: <N>} (\texttt{input\_read.f90}); an
\texttt{input.inp} without the key keeps 500 and reproduces its earlier
result. The grid-sized module arrays became allocatable (lower bound
\texttt{1-Ng} unchanged) and are allocated by
\texttt{allocate\_grid\_arrays} once parsing has resolved the key, so
nothing below is affected other than that $N$ may now differ from 500.
Values under 10 are refused. Two consequences appear later: the base
refinement pairs of \S\ref{sec:basegrid}, and the restart record count
of \S\ref{sec:loadic}.

The first interior cell $j=1$ sits at the lower boundary radius
$r = 1$ (lengths are normalized to the planetary radius $R_0$); the
last interior cell $j=N$ sits at $r_{\max}$, which is either the
Roche/Hill truncation $r_{\max} = (3\,M_*/M_p)^{-1/3}\,\tilde a$ in
Roche mode or a user-chosen outer radius in spherical mode
(\texttt{input\_read.f90}).

All boundary conditions are applied to the \emph{ghost cells only};
the interior cells $1,\dots,N$ are never overwritten. The fluxes at
the two boundary interfaces then emerge from the same HLLC Riemann
solver used everywhere else, with the ghost states as the outside
input. In this sense every BC in the code is a \emph{weak} (flux-mediated)
condition: even the ``Dirichlet'' base values act on the interior only
through the Riemann problem at the interface $j = 1/2$.

\subsection{The three entry points (\texttt{Apply\_BC.f90})}

All BC logic lives in one module,
\texttt{src/modules/states/Apply\_BC.f90}:

\begin{itemize}
\item \texttt{Apply\_BC(u,u)} --- conservative-variable wrapper:
  converts $U \to W$ (primitive), calls \texttt{Apply\_BC\_W}, converts
  back. Called after every Runge--Kutta stage in the main loop
  (\texttt{EXHALE\_main.f90}), before every residual evaluation, and
  by the steady-state Newton solver each time it unpacks a trial state.
\item \texttt{Apply\_BC\_W(W,W)} --- the actual assignment of the four
  ghost states: lower ghosts via \texttt{BC\_component\_constrho}
  (\S\ref{sec:lower}), upper ghosts via extrapolation
  (\S\ref{sec:upper}).
\item \texttt{Rec\_BC(WL,WR,...)} --- the same conditions applied to the
  \emph{reconstructed} left/right interface states after PLM or WENO3
  reconstruction (\texttt{Reconstruction.f90}), so that the boundary
  stencils of the higher-order scheme see BC-consistent data rather
  than raw extrapolated slopes.
\end{itemize}

In the Newton steady-state solver the ghost cells are deliberately
\emph{not} unknowns: the residual is defined on the interior state
only, and \texttt{Apply\_BC} regenerates the ghosts from the interior
at every function evaluation
(\texttt{src/modules/time\_step/steady\_newton.f90}). The BCs are
therefore identical between the time-marching and Newton phases by
construction.

\subsection{Normalization constants tied to the base}
\label{sec:normalization}

The lower BC is intertwined with the code-wide normalization
(\texttt{input\_read.f90}):
\begin{equation}
\mu = \frac{1 + 4\,Y}{1 + Y}\, m_{\rm H}, \qquad
v_0 = \sqrt{\frac{k_B T_0}{m_{\rm H}}}, \qquad
p_0 = n_0\, m_{\rm H} v_0^2 = n_0 k_B T_0, \qquad
b_0 = \frac{G M_p\, m_{\rm H}}{k_B T_0 R_0},
\end{equation}
where $Y \equiv n_{\rm He}/n_{\rm H}$ is the \texttt{He/H number
ratio} from \texttt{input.inp}, $n_0$ is the user-specified total
\emph{nuclei} number density at the base, and $T_0$ the base
temperature. (The code variable \texttt{mu} is $m_{\rm H}$ itself;
the composition factor appears separately as \texttt{rho\_bc}, see
below.) Densities are stored as mass densities in units of
$n_0 m_{\rm H}$, number densities in units of $n_0$, pressures in
units of $p_0$, temperatures in units of $T_0$, velocities in units of
the isothermal hydrogen sound speed $v_0$.

The equation-of-state closure used to convert between the hydrodynamic
pressure and the temperature handed to the ionization/radiation module
is the dimensionless ideal-gas law including free electrons
(\texttt{composition.f90}):
\begin{equation}
T = \frac{p}{n_{\rm tot} + n_e},
\label{eq:eos}
\end{equation}
with $n_{\rm tot}$ the total \emph{atomic} (nuclei) number density and
$n_e$ from \texttt{calc\_ne}. Under the default \texttt{eos\_metals 1}
policy (2026-06-13) both include the trace metals --- their nuclei in
$n_{\rm tot}$ and their electrons in $n_e$ --- consistently with the
charge balance inside the ionization solver; \texttt{eos\_metals 0} in
\texttt{metals.inp} restores the legacy H/He-only budget
(\S\ref{sec:metals-bc}).

\section{Lower boundary: fixed reservoir with a one-way valve}
\label{sec:lower}

The two lower ghost cells are set by
\texttt{BC\_component\_constrho} (\texttt{Apply\_BC.f90}, lines
31--80). In the cell-index-first notation of this memo
(\texttt{W(cell,component)}; the code stores the transpose):
\begin{lstlisting}[style=code]
W(ghost,1) = rho_bc                      ! density: Dirichlet
if (base_v_massflux .and. base_flux_const > 0) then    ! mass-flux base
   W(ghost,2) = base_flux_const/(rho_bc*r(ghost)**2)
else if (valve_eps > 0) then                           ! smooth valve
   W(ghost,2) = 0.5*( W(1,2) + sqrt(W(1,2)**2 + valve_eps**2) )
else                                                   ! legacy hard valve
   W(ghost,2) = max( W(1,2), 0.0 )
endif
if (hydrostatic_base) then                             ! dp/dr continuous
   W(ghost,3) = W(1,3) + (W(2,3) - W(1,3))/(r(2) - r(1))*(r(ghost) - r(1))
else if (base_ghost_T_continuous) then                 ! dT/dr = 0
   W(ghost,3) = (ntot_bc + dp_bc)*W(1,3)/n_part_cell1
else                                                   ! T = T0 (default)
   W(ghost,3) = ntot_bc + dp_bc
endif
\end{lstlisting}
The default closures (legacy valve, isothermal ghost pressure) give both
ghost cells the \emph{same} state, so there is no gradient across them;
the mass-flux velocity and the hydrostatic ghost pressure are evaluated at
$r(\mathrm{ghost})$ and therefore do differ between $j=0$ and $j=-1$. The
three components implement three distinct ideas, discussed in turn; the
opt-in velocity and pressure branches are collected in
\S\ref{sec:stabilizers}.

\subsection{Density: a fixed-composition reservoir}

The dimensionless ghost density is the constant
\begin{equation}
\rho_{\rm bc} = \frac{m_{\rm gas}/{\rm H}}{1 + Y},
\qquad
m_{\rm gas}/{\rm H} = 1 + 4Y + \sum_X a_X A_X,
\qquad a_X \equiv \frac{n_X}{n_{\rm H}},
\label{eq:rhobc}
\end{equation}
the mean gas mass per H/He nucleus in units of $m_{\rm H}$
(\texttt{comp\_rho\_bc()} $=$ \texttt{comp\_mass\_per\_H()}$/(1{+}Y)$ in
the code). Because mass densities
are normalized by $n_0 m_{\rm H}$, this pins the ghost cells to
``H/He nuclei density $= n_0$ with the prescribed composition'': the
hydrogen nuclei density at the base is $n_{\rm H} = n_0/(1+Y)$, the
helium density $Y$ times that, and each metal $a_X n_{\rm H}$. The
base therefore behaves as an \emph{infinite isochoric reservoir}:
however much mass the wind removes, the boundary cell always refills
to the same density.

The metal mass term $\sum_X a_X A_X \approx 0.016$ at solar abundances
(a ${\sim}1.2\%$ effect) is included by the default
\texttt{eos\_metals 1} policy, which also adds the metal mass to
\texttt{calc\_rho}; \texttt{eos\_metals 0} restores the legacy
H/He-only trace approximation, $\rho_{\rm bc} = (1+4Y)/(1+Y)$. That
legacy expression is therefore \emph{only} the value of a metals-off
atomic base; it is not the general form.

\paragraph{One source for the base composition.}
$\rho_{\rm bc}$, $m_{\rm gas}/{\rm H}$ (\texttt{mass\_per\_H}) and
$n_{\rm tot,bc}$ (\texttt{ntot\_bc}) are all produced by the three
functions \texttt{comp\_rho\_bc}, \texttt{comp\_mass\_per\_H} and
\texttt{comp\_ntot\_bc} of
\texttt{src/modules/functions/composition.f90}, which
\texttt{input\_read} calls once; that module is the single place where
the base composition policy lives, so no code path can disagree with
another about what the base is made of. Three switches move the numbers
inside it:
\begin{itemize}
\item \texttt{eos\_metals} (\texttt{metals.inp}; default 1 since
  2026-06-13) puts the trace-metal mass into $m_{\rm gas}/{\rm H}$ and
  the trace-metal nuclei into $n_{\rm tot,bc} = (1 + Y + \sum_X
  a_X)/(1+Y)$. With \texttt{eos\_metals 0}, or with no metals, both
  collapse to the H/He values $1+4Y$ and $1$.
\item \texttt{Molecular base: True} (Tier-2a; auto-enabled whenever the
  molecular chemistry is on) subtracts from $n_{\rm tot,bc}$ the
  particles lost to the H nuclei bound into H$_2$ at the base: two H
  nuclei make one molecule, so a bound H fraction $x_2$ costs $x_2/2$
  particles (\texttt{h2\_bound\_fraction}). This changes the particle
  count only --- $\rho_{\rm bc}$, hence the mass reservoir, is untouched,
  and the species arrays stay atomic.
\item A \texttt{base.inp} handoff supplies $q_{\rm H_2}$ at the base
  (\texttt{q\_H2\_base}) and then takes precedence over the
  chemical-equilibrium fit that \texttt{h2\_mixing\_ratio\_base}
  otherwise evaluates at $(p_{\rm base}, T_0)$; \texttt{input\_read}
  echoes which of the two set $n_{\rm tot,bc}$. Rationale:
  \texttt{docs/base\_composition\_handoff\_plan.md}.
\end{itemize}

\subsection{Pressure: dynamically updated to pin the base temperature}
\label{sec:dpbc}

The ghost pressure is $p_{\rm gh} = n_{\rm tot,bc} + \delta p_{\rm bc}$,
where $n_{\rm tot,bc} = (1 + Y + \sum_X a_X)/(1+Y)$ is the (constant)
total nuclei density at the base in units of $n_0$
(\texttt{ntot\_bc}; $=1$ for H/He only, and reduced by the H$_2$-bound
share under a molecular base) and $\delta p_{\rm bc}$ starts
at the placeholder $10^{-10}$ and is recomputed after every
ionization-equilibrium solve as the base electron density
(\texttt{ionization\_equilibrium.f90}):
\begin{equation}
\delta p_{\rm bc}
   = \frac{n_{\rm HII} + n_{\rm HeII} + 2\,n_{\rm HeIII}
           + \sum_i q_i\, n_i^{\rm met}}{n_0}
     \bigg|_{j\,=\,1-N_g}
   = \frac{n_e(\mbox{base})}{n_0},
\end{equation}
with $q_i$ the metal ion charges (included under
\texttt{eos\_metals 1}, dropped under \texttt{eos\_metals 0}).
The point of this correction is visible through the EOS closure
(\ref{eq:eos}): at the ghost cells $n_{\rm tot} = n_{\rm tot,bc}$ (by
the density BC) and $n_e = \delta p_{\rm bc}$, so
\begin{equation}
T_{\rm gh} \;=\; \frac{n_{\rm tot,bc} + \delta p_{\rm bc}}
                      {n_{\rm tot,bc} + \delta p_{\rm bc}}
          \;=\; 1
  \qquad\Longleftrightarrow\qquad T = T_0 .
\end{equation}
In other words, the pair (fixed $\rho$, dynamically corrected $p$)
is really the pair \emph{(fixed nuclei density $n_0$, fixed
temperature $T_0$)}: as photoionization develops at the base and free
electrons add partial pressure, the ghost pressure is raised by
exactly the electron share so that the base stays isothermal at
$T_0$ instead of being artificially cooled. For a near-neutral base
(the usual situation, since the base optical depth is huge)
$\delta p_{\rm bc}$ is tiny and the BC reduces to $p = p_0$.

The metal-electron share of $\delta p_{\rm bc}$ is small (even with
all low-ionization-potential metals singly ionized at the base it is
${\lesssim}4\times10^{-4}$), but it is carried anyway under
\texttt{eos\_metals 1} for consistency with \texttt{calc\_ne}.

\subsection{Velocity: the one-way valve}
\label{sec:valve}

The ghost velocity copies the first interior cell when the gas moves
outward and is clamped at zero when it moves inward:
\begin{equation}
v_{\rm gh} = \max(v_1,\, 0).
\end{equation}
Neither a wall ($v_{\rm gh}=0$) nor a free condition
($v_{\rm gh}=v_1$) would be physical here:
\begin{itemize}
\item a \emph{wall} would starve the wind --- the escaping flow is
  launched by photoionization heating \emph{above} the base, and the
  mass it carries must be resupplied through the lower boundary;
\item a \emph{free} condition combined with the fixed-density ghosts
  would let the reservoir act as an unphysical mass sink whenever the
  interior gas sags downward, swallowing arbitrary amounts of gas at
  fixed $\rho_{\rm bc}$.
\end{itemize}
The valve resolves the asymmetry: \emph{outflow is allowed (the
reservoir feeds the wind), sustained inflow is refused (the reservoir
does not absorb a backflow)}. Physically the boundary models the top
of an essentially infinite, hydrostatic lower atmosphere held at
$(n_0, T_0)$.

\paragraph{Smooth valve (EXHALE extension).}
The hard $\max$ has a kink at $v_1 = 0$ --- and, as discussed in
\S\ref{sec:breathing}, $v_1 = 0$ is precisely where the base
dynamics lives. The kink is invisible to the explicit time-marching
scheme but fatal to the steady-state Newton solver, whose line search
cannot cross a non-differentiable point of the residual. The input
key \texttt{Valve eps:} therefore enables the smooth approximation
\begin{equation}
v_{\rm gh} = \tfrac12\!\left( v_1 + \sqrt{v_1^2 + \varepsilon^2} \right),
\end{equation}
which is $C^\infty$, tends to $v_1$ for $v_1 \gg \varepsilon$, tends
to $0$ for $v_1 \ll -\varepsilon$, and biases the valve by at most
$+\varepsilon/2$ near $v_1 \sim 0$. With
$\varepsilon = 0$ (default) the legacy hard valve is used.

\subsection{The ``breathing base''}
\label{sec:breathing}

Because the base velocity is a free variable governed by an
\emph{asymmetric} boundary (an open pipe outward, a wall inward), the
near-base flow does not settle onto the exact $v = 0$ fixed point but
onto a small limit cycle around it. One cycle proceeds as follows:

\begin{enumerate}
\item \textbf{Exhale} ($v_1 > 0$): the wind draws mass up; the
  Riemann problem against the fixed-$\rho_{\rm bc}$ ghosts supplies
  gas at base density --- the boundary acts as a mass source.
\item \textbf{Overshoot}: if momentarily more mass is lifted than the
  heating profile above can carry, the over-supplied column begins to
  sink back under gravity; $v_1$ turns negative.
\item \textbf{Inhale} ($v_1 < 0$): the valve closes
  ($v_{\rm gh} = 0$). The descending gas piles up on the closed
  boundary, raising $\rho$ and $p$ in the first few interior cells ---
  a transient wall reflection.
\item \textbf{Rebound}: the accumulated pressure pushes the column
  back out; $v_1$ turns positive and the cycle repeats.
\end{enumerate}

This inhale/exhale oscillation is the \emph{breathing base}. It is an
intrinsic consequence of the valve-type BC --- present in the original
ATES as well --- and not a numerical defect. Two regimes should be
distinguished (see \texttt{docs/initial\_condition\_benchmark.pdf},
\S7):

\begin{description}
\item[Self-limiting (normal).] In the time-marching hydrodynamics the
  pressure built up during the inhale phase immediately produces the
  restoring force of step 4, so the amplitude saturates. The
  oscillation stays confined to the first few cells
  ($r \lesssim 1.1\,R_p$) and has no measurable effect on the wind
  region or on $\dot M$.
\item[Self-amplifying (pathological).] If the \emph{global} initial
  condition is inconsistent (the canonical example: a naive full-Parker
  density profile, whose hot column drains onto the base), each cycle
  deposits more mass than the previous one and the breathing grows
  until the run diverges. The cure is a consistent IC (hydrostatic
  density, optionally with a warm seed overlay), not a change of the
  BC.
\end{description}

\paragraph{Numerical consequences.}
Harmless as it is physically, the breathing base interacts with three
parts of the code, each of which carries a dedicated counter-measure:
\begin{itemize}
\item \emph{Convergence monitoring.} The breathing cells never stop
  oscillating, so a convergence criterion of the form
  $\max_j|\Delta u|$ over the whole grid can stall at the
  $10^{-2}$--$10^{-1}$ level indefinitely. ATES therefore evaluates
  its momentum-flatness and $du$/$dtu$ diagnostics only over the wind
  region $[\,j_{\min} : N\,]$, where $j_{\min}$ is the first cell with
  $r \ge r_{\rm esc}$ (the \texttt{Escape radius} input;
  \S\ref{sec:resc}). Cases with unusually violent breathing
  (HD~189733~b is the standing example) can still fail the strict
  $du$ threshold; reducing CFL, modifying the valve, or replacing the
  boundary Riemann solver did \emph{not} remove the oscillation in our
  experiments --- it is robustly attached to the BC type. A periodic
  Shapiro (1970) low-pass filter, the stabilizer the reference code
  CETIMB (Koskinen et al.\ 2013a) applies to exactly this
  gravity-unbalanced sound-wave instability, \emph{does} damp the
  oscillatory component (\S\ref{sec:stabilizers}), but it biases the
  momentum and pushes the clean cold-IC cases off the wind branch into
  infall, so it is kept opt-in rather than a default.
\item \emph{Newton solver.} The $v_1 = 0$ kink of the hard valve sits
  exactly on the breathing orbit; the smooth valve of
  \S\ref{sec:valve} exists for this reason.
\item \emph{Advection post-processing.} The \texttt{\_adv} correction
  integrates ionization states along streamlines; during the inhale
  phase ($v \le 0$) the streamline of a base cell exits the domain
  through the lower boundary and the correction is ill-defined (this
  produced the historical temperature spike/sawtooth at
  $r \approx 1.1$). The post-processor therefore keeps the
  equilibrium solution and skips the advection correction wherever
  $v \le 0$ at the base.
\end{itemize}

\subsection{Base grid resolution (2026-08-11)}
\label{sec:basegrid}

Not every alternating pattern at the base is the breathing limit cycle.
On the \texttt{Mixed} grid the base region is $N_{\rm low}$ uniform cells
of width $\Delta r_{\rm base}$ stacked on $r=1$, and the physical
requirement on that width is that it resolve the base density scale
height $H = kT/(\mu g)$. Where $H/\Delta r_{\rm base}$ is of order a few
cells, the discretization supports a \emph{stationary} $2\Delta r$
entropy (contact) mode that nothing in the scheme damps --- HLLC resolves
a contact of zero speed exactly, the gravity source is cell-local, and
the WENO3 pressure gradient sees only interface pressures. The measured
dependence is an alternating amplitude in $\ln\rho$ of ${\sim}10^{-2}$
for $H/\Delta r < 5$ against ${\sim}10^{-4}$ for $H/\Delta r > 100$
(\texttt{docs/hd189\_base\_checkerboard.md}). High-gravity planets have
the least margin: on the default grid \texttt{write\_setup\_report}
echoes 26.9 cells per $H$ at $T_{\rm eq}$ for HD\,189733\,b against
102.3 for WASP-121\,b.

The resolution is therefore an input key,
\texttt{Base grid [dr,cells]: <dr\_base> [<N\_low\_cells>]}
(\texttt{input\_read.f90}; defaults \texttt{2.0e-4 50}, ignored by the
\texttt{Uniform} and \texttt{Stretched} grid types). The two numbers are
not independent --- their product is the radial extent of the uniform
region, $0.01\,R_p$ with the defaults --- so refining the base at fixed
extent means dividing $\Delta r_{\rm base}$ and multiplying $N_{\rm low}$
by the same factor. Because those extra base cells are taken out of the
same total $N$, the runtime \texttt{Grid cells:} key of
\S\ref{sec:overview} is what lets a refinement keep the outer stretched
region unchanged; the paired settings are
\begin{center}\small
\begin{tabular}{@{}ll@{}}
\toprule
\texttt{Base grid [dr,cells]:} & \texttt{Grid cells:} \\
\midrule
\texttt{1.0e-4 100}  & 500 \\
\texttt{5.0e-5 200}  & 658 \\
\texttt{2.5e-5 400}  & 916 \\
\texttt{1.25e-5 800} & 1375 \\
\bottomrule
\end{tabular}
\end{center}
Note that spelling the default out explicitly
(\texttt{Base grid [dr,cells]: 2.0e-4 50}) does \emph{not} reproduce a
run made without the key: the compiled-in default is a default-real
literal, $2.5\times10^{-8}$ relative below the exact double the
list-directed read produces, which moves every base cell by that amount.
The full measurement --- what sets the mode, which remedies were tried,
and what each grid step bought --- is
\texttt{docs/hd189\_base\_checkerboard.md}.

\subsection{Optional stabilizers and the selectable base mode (EXHALE)}
\label{sec:stabilizers}

EXHALE adds six opt-in controls of the base behavior. The first three were
ported while studying the breathing/infall behavior against CETIMB (Koskinen
et al.\ 2013a); the Shapiro
filter in particular was adopted after Huang et al.\ (2023) --- whose escape
model is built on CETIMB --- pointed to CETIMB's use of exactly this low-pass to
damp the gravity-unbalanced base sound waves; that reference is what suggested
trying it here. The last three (2026-08) came out of the HD\,189733\,b base
checkerboard and the HD\,209458\,b stagnant-layer investigations. All six
\emph{default to the legacy behavior} --- the numeric ones are off at
$\varepsilon \le 0$, the ones selected by name default to \texttt{valve},
\texttt{density} and \texttt{isothermal} --- so a run that does not ask for
them is unchanged. The measured effects of the first three are
tabulated in \S\ref{sec:bc-matrix}.

Two of them --- \texttt{Hydrostatic base} and \texttt{Base ghost
temperature: continuous} --- set the \emph{same} quantity, the ghost
pressure, and are therefore mutually exclusive: \texttt{Hydrostatic base}
is tested first in \texttt{BC\_component\_constrho} and
\texttt{input\_read} warns when both are given.

\begin{description}
\item[Shapiro filter (\texttt{Shapiro filter: <eps> [<every>]}).] Every
  \texttt{every} steps a weak two-step Shapiro (1970) 1--2--1 low-pass filter
  $u_j \mathrel{+}= \tfrac{\varepsilon}{4}(u_{j-1} - 2u_j + u_{j+1})$ is applied
  to the conservative variables. This is CETIMB's cure for the
  gravity-unbalanced sound waves and it does smooth the breathing transient.
  \emph{However}, on the otherwise-clean cold-IC cases it injects a systematic
  momentum bias that drives the solution off the transonic-wind saddle into the
  infall branch ($v<0$ almost everywhere; see the \texttt{\_Shap} rows of
  Table~\ref{tab:bc-dens}). It is therefore opt-in --- useful only for cases
  that genuinely breathe, not as a global default.
\item[Mass-flux base velocity (\texttt{Base velocity: massflux}).] Instead of
  the one-way valve (\S\ref{sec:valve}), the base velocity is set from the
  steady mass-flux continuity $v_0 = F_c/(\rho_0 r_0^2)$ with $F_c$ a slow
  exponential moving average of $\rho v r^2$ over the wind region. It is
  physically benign --- it never causes infall on its own and gives the cleanest
  base (it suppresses the $\pm$m\,s$^{-1}$ valve oscillation) --- but it slows
  convergence ($\sim$5$\times$ more steps), so the fast valve stays the default.
\item[Pressure-anchored base (\texttt{Base BC: pressure [<p\_ubar>]}).] Instead
  of fixing the base \emph{number density} at $n_0$ (\S\ref{sec:lower}), this
  fixes the base \emph{pressure} at \texttt{p\_ubar} microbar (default $1\,\mu$bar,
  the CETIMB lower boundary) and derives $n_0$ from $p = n_0 k_B T_0\,
  n_{\rm tot,bc}$. Because the isothermal base is over-determined (it already
  pins $T = T_0$), this is a re-parameterization of the same ghost; its physical
  effect is simply a much less dense base ($1\,\mu$bar $\Rightarrow
  n_0 \approx 5\times10^{12}$, $\sim$20$\times$ below $10^{14}$), which shrinks
  the dense-base gravity source $\rho g$ that feeds the breathing. Empirically a
  density-anchored run at the same low $n_0$ ($\log n_0 = 12.5$) gives an almost
  identical converged wind, confirming that the lever is the base \emph{density},
  not the anchoring variable (Tables~\ref{tab:bc-pres} and~\ref{tab:bc-dens}).
\item[Base ghost temperature (\texttt{Base ghost temperature:
  isothermal|continuous}, 2026-08-11).] Selects the temperature closure of
  the lower ghost cells; the density anchor $\rho_{\rm bc}$ is untouched
  either way. \texttt{isothermal} (default, legacy) is the pin of
  \S\ref{sec:dpbc}: the ghost pressure is held at
  $n_{\rm tot,bc} + \delta p_{\rm bc}$, i.e.\ $T_{\rm gh} = T_0$.
  \texttt{continuous} imposes $\mathrm{d}T/\mathrm{d}r = 0$ at the base
  face instead --- the ghost keeps the base composition (the same
  $n_{\rm tot,bc} + \delta p_{\rm bc}$ particles at $\rho_{\rm bc}$) but
  carries the temperature of the first interior cell,
  \begin{equation}
  p_{\rm gh} = (n_{\rm tot,bc} + \delta p_{\rm bc})\,T_1 ,
  \qquad
  T_1 = \frac{p_1}{(n_{\rm tot} + n_e)_1},
  \end{equation}
  with $(n_{\rm tot}+n_e)_1$ (\texttt{n\_part\_cell1}) refreshed by the
  composition solve, so the ghost pressure stays a differentiable function
  of the interior pressure --- what the JFNK line search needs --- while
  the ionization state it divides by is lagged like every other
  composition quantity in a hydro step
  (\texttt{Apply\_BC.f90}, lines 55--77; \texttt{input\_read.f90}, lines
  588--602). The motivation is that the $T_0$ pin has no physical backing
  --- the radiative equilibrium of the lower atmosphere is outside the
  model --- and becomes a contact discontinuity sitting on the boundary
  when the first cell cools well below $T_0$. Measured on the production
  HD\,189733\,b configuration
  (\texttt{docs/hd189\_base\_checkerboard.md}, \S13): $T_{\rm gh}$ goes
  from $1183$~K to $540$~K against $T_1 \approx 494 \to 540$~K,
  $\rho_1/\rho_{\rm gh}$ from 2.33 to 0.95, and the alternating $\ln\rho$
  amplitude over cells 1--12 from $2.6\times10^{-2}$ to
  $9.2\times10^{-3}$, at an unchanged $\log_{10}\dot M$ (9.04 vs.\ 9.05).
  Note that with \texttt{Base BC: pressure} the microbar target is imposed
  at $T_0$ when $n_0$ is derived, so under a continuous-$T$ ghost only the
  base \emph{density} remains anchored; \texttt{input\_read} says so.
\item[Hydrostatic base (\texttt{Hydrostatic base: True}).] The legacy ghost
  pressure is a constant across the two ghost cells, so
  $\mathrm{d}p/\mathrm{d}r$ is flattened to ${\approx}0$ at the base ---
  while steady momentum balance there needs
  $\mathrm{d}p/\mathrm{d}r = -\rho g$. This flag replaces the ghost pressure
  by the linear extrapolation of the \emph{interior} pressure gradient
  (cells 1 and 2),
  \begin{equation}
  p_{\rm gh}(r) = p_1 + \frac{p_2 - p_1}{r_2 - r_1}\,(r - r_1),
  \end{equation}
  so that $\mathrm{d}p/\mathrm{d}r$ is continuous through the base face and
  can balance gravity there --- the base-face momentum residual being the
  source of the breathing (\S\ref{sec:breathing}), which for strongly bound
  planets (large $b_0$, e.g.\ HD\,189733\,b) is $O(10$--$100)$. The density
  stays anchored at $\rho_{\rm bc}$; the base temperature then floats
  slightly off $T_0$, since neither $T$ nor $p$ is fixed any more
  (\texttt{Apply\_BC.f90}, lines 55--72).
\item[Low-Mach contact-mode damping (\texttt{Low-Mach damping: <eps4>
  [<M\_th>]}, 2026-08-13).] Not a boundary condition, but aimed at the same
  near-base region. The HLLC flux carries the contact/entropy wave at
  $S_* \sim v$, so the dissipation it applies to that family vanishes as
  the flow stagnates and a $2\Delta r$ mode in $v$ and $T$ becomes
  marginally damped --- which is what a shell where metal cooling has
  stopped the flow looks like. The key adds a gated fourth-difference
  (Jameson--Schmidt--Turkel) stress to the numerical \emph{momentum} flux,
  together with its work term $v\,D_p$ in the energy flux so the pair
  conserves total energy; the mass flux is untouched. The gate closes the
  term identically above $M > M_{\rm th}$ (default $10^{-3}$). It lives in
  \texttt{src/modules/flux/low\_mach\_dissipation.f90} and is applied
  inside \texttt{RK\_rhs}, the one routine the steady residual assembly
  also calls, so the marching loop and the Newton residual see the
  \emph{same} equation --- unlike the Shapiro filter, which touches only
  the marching state, leaving Newton to solve a system whose mode the
  filter had been hiding. Off by default ($\varepsilon_4 \le 0$ skips the
  code path entirely). \emph{Caveats, both measured}
  (\texttt{docs/hd209\_metal\_stagnation.md}, \S\S8--9): the coefficient
  has an operating window --- \texttt{input\_read} warns above the explicit
  stability bound $\varepsilon_4 < 1/(16\,\mathrm{CFL})$, and on the
  \texttt{examples/16} warm restart only $5\times10^{-3}$ and $10^{-2}$
  gave a clean Newton finish --- and on the \emph{cold} start of that same
  configuration no coefficient helped at all: off, $5\times10^{-3}$,
  $10^{-2}$ and $2\times10^{-2}$ all bottom out at the same
  ${\sim}10^{-3}$ residual floor, so the cold hand-off state appears to lie
  outside the Newton basin and the damping is not the lever there.
\end{description}

\paragraph{Implementation note.}
The cell-by-cell ionization-equilibrium sweep is now OpenMP-parallel over cells
(the metal and charge-exchange coefficients in each cell are kept thread-private);
the parallel result is bit-identical to the serial one, so none of the BC/IC
behavior below is affected. See \texttt{docs/openmp\_parallelization.md}.

\subsection{Temperature, ionization, and metals at the lower boundary}
\label{sec:metals-bc}

No boundary condition is imposed on the temperature or on any
ionization fraction. The ghost cells are included as ordinary cells in
the ionization-equilibrium sweep
(\texttt{ionization\_equilibrium.f90} loops over
$j = N{+}N_g,\dots,1{-}N_g$), so their $T$, H/He stages, He~$2^3$S,
and all metal stages come out of the same coupled MINPACK solve as the
interior, evaluated at the ghost state $(\rho_{\rm bc},\,
n_{\rm tot,bc}{+}\delta p_{\rm bc})$. Combined with \S\ref{sec:dpbc} this makes the
effective thermodynamic boundary condition ``$n = n_0$, $T = T_0$,
ionization in local equilibrium''. Metal abundances enter only through
the runtime \texttt{metals.inp} fractions; the boundary needs no
special metal treatment: their abundances ride on the (fixed-composition)
density BC, and under the default \texttt{eos\_metals 1} policy their
mass, electrons, and nuclei are carried consistently through
\texttt{calc\_rho}/\texttt{calc\_ne}/\texttt{calc\_ntot}
(\texttt{eos\_metals 0} restores the legacy trace approximation in
which they never appear in the mass or momentum budget).

\section{Upper boundary: supersonic free outflow}
\label{sec:upper}

The upper ghost cells are filled by extrapolation
(\texttt{Apply\_BC\_W}, lines 95--99):
\begin{lstlisting}[style=code]
do k = 1, Ng
   W(N+k,:) = W(N,:)                              ! PLM / first order
   if (use_weno3) W(N+k,:) = 2*W(N+k-1,:) - W(N+k-2,:)  ! WENO3
enddo
\end{lstlisting}
i.e.\ zeroth-order (copy) extrapolation of all primitive variables for
the PLM scheme, upgraded to linear extrapolation when WENO3
reconstruction is active (the wider stencil of WENO3 is sensitive to
the artificial flattening a plain copy would imply). \texttt{Rec\_BC}
applies the same rule to the reconstructed interface states.

The justification is characteristic-based: by the time the flow
reaches $r_{\max}$ the wind is expected to be supersonic, so all three
characteristics leave the domain and no physical information may be
prescribed from outside; extrapolation is then the consistent
``transparent'' closure. Two caveats follow directly:
\begin{itemize}
\item If the flow is \emph{not} supersonic at $r_{\max}$ (observed for
  deep Roche-lobe-overflow geometries, e.g.\ WASP-121b Case~D, which
  is subsonic at L1), one in-going characteristic formally remains and
  the extrapolation BC is no longer strictly well-posed; in practice
  the scheme remains stable but the outer solution should be treated
  with caution (geometry-limited regime).
\item The BC carries no information about the medium beyond $r_{\max}$
  (stellar wind, confinement pressure); EXHALE models the planetary
  outflow only.
\end{itemize}

No ionization or temperature condition is applied at the top either;
the ghost cells participate in the equilibrium sweep like all others,
with their extrapolated $(\rho, p)$.

\section{Boundary-adjacent settings}
\label{sec:adjacent}

These are not boundary conditions of the hydrodynamic update, but they
are boundary choices of the model and are easily confused with BCs.

\subsection{Radiative transfer: irradiation from the top}

The stellar EUV/X-ray flux enters through the \emph{outer} boundary.
The column densities that attenuate the ionizing flux are integrated
top-down (\texttt{calc\_column\_dens}, \texttt{utilities.f90}): the
outermost ghost cell is assigned its own local column
$N(N{+}N_g) = n\,\Delta r$, and the recursion
$N(j) = N(j{+}1) + n_j\,\Delta r_j$ proceeds inward, i.e.\ the optical
depth is (essentially) zero at the top of the grid and grows toward
the base. This is the radiative-transfer analogue of a boundary
condition: the unattenuated stellar spectrum is imposed at $r_{\max}$,
and the base receives whatever survives the overlying column. Because
the base column is huge, the lower boundary is effectively in
radiative shadow --- consistent with holding it at $(n_0, T_0)$.

\subsection{Escape radius and the convergence window}
\label{sec:resc}

The \texttt{Escape radius [R\_p]:} input ($r_{\rm esc}$, default 1.5)
defines $j_{\min}$, the first cell with $r \ge r_{\rm esc}$
(\texttt{define\_grid.f90}). All steady-state diagnostics ---
momentum-flux flatness $\rho v r^2 = \mathrm{const}$, $du$, $dtu$ ---
are evaluated over $[\,j_{\min} : N\,]$ only, deliberately excluding
the breathing base (\S\ref{sec:breathing}). A guard added for deep-RLOF
geometries (EXHALE): if $r_{\rm esc} \ge r_{\max}$ the window would be
empty (and \texttt{maxval}/\texttt{minval} over an empty section
return $\mp$\texttt{HUGE}, producing a spurious instant
``convergence''); in that case the window is clamped to the outer half
of the domain and a loud warning asks the user to lower
$r_{\rm esc}$.

\paragraph{2026-08-15: the layer below the window is measured, not accepted
silently.} The acceptance window is still $[\,j_{\min}:N\,]$, but the layer
$[\,1:j_{\min}{-}1\,]$ excluded from it is no longer unreported.
\texttt{resid\_relnorm\_below\_escape}
(\texttt{src/modules/time\_step/steady\_newton.f90}) applies the same
volume-weighted relative norm to that layer and locates the cell carrying
the largest cell-wise residual in it;
\texttt{write\_resid\_below\_escape} prints one line at the start and end
of the PTC solve and at the start and end of the JFNK solve. The value is
printed and never tested --- nothing in the solve depends on it.

The momentum row of that diagnostic is \emph{not} scaled by
$|u_2| = |\rho v|$. The layer is quasi-hydrostatic ($|v|\sim10$\,cm\,s$^{-1}$),
so that denominator collapses and the ratio is large wherever the flow is
slow, whether or not the state is right: converged states read 0.2--10 on
it, the metals-off solution nobody doubts reading the largest. The physical
scale of the momentum equation there is the hydrostatic pair
$|\mathrm{d}p/\mathrm{d}r| \sim \rho g$, so the momentum row is divided by
the gravitational force density
\begin{equation}
s_{\rm grav}(j) = |\rho_j|\,
   \frac{\big|G\phi_i(j) - G\phi_i(j{-}1)\big|}{\Delta r_j},
\end{equation}
the same discrete potential difference the momentum source term uses
(\texttt{Source.f90}); the printed number is then the fractional violation
of hydrostatic balance in the layer. On that scale the converged
HD\,209458\,b solutions --- including the metals-off one that read 9.99
before --- all hold the layer at $10^{-5}$ to $1.4\times10^{-4}$ of
$\rho g$, comparable to the wind-region \texttt{Resid tol}. The mass and
energy rows keep the $|u|$ scale and print separately.

Widening the acceptance window to include the layer was considered and is
\emph{rejected by measurement}: on the old $|u|$ scale a converged state
carries $0.2$--$10$ there, four to six orders above \texttt{Resid tol}, so
the widened window would reject solutions that are correct. Both
measurements: \texttt{docs/hd209\_metal\_stagnation.md}, \S\S9.4 and 10.

\subsection{Domain truncation (where the upper boundary sits)}

In Roche mode the grid is truncated at the Hill radius
$r_{\max} = (3 M_*/M_p)^{-1/3}\tilde a \approx$ L1; in spherical mode
at the user-supplied \texttt{Outer radius [R\_p]:}. The choice changes
where the free-outflow BC is applied, and hence how defensible the
supersonic assumption is (\S\ref{sec:upper}); it does not change the
BC itself.

\section{Initial conditions (\texttt{set\_IC.f90})}
\label{sec:ic}

\paragraph{Three concepts to fix first: \textbf{cold IC}, \textbf{warm IC}, and
\textbf{starter}.} For the reader's convenience, every IC family below is a
variation on three ideas. The distinction matters because it controls whether
the run starts \emph{inside} the wind basin or has to be pushed there.
\begin{description}
\item[\textbf{Cold IC.}] The default. A fully \emph{neutral}, isothermal
  ($T = T_0$) \emph{hydrostatic} atmosphere with only a small outward velocity
  seed. Nothing is pre-heated or pre-ionized: the simulation grows the
  temperature, ionization, and wind from scratch. Because it is built from
  EXHALE's own base BC (fixed $\rho_{\rm bc}$, $T_0$), it is automatically
  \emph{consistent with the lower boundary} and relaxes cleanly to the
  transonic wind.
\item[\textbf{Warm IC} (warm seed / warm start).] An initial state that already
  carries a \emph{heated, ionized, flowing} profile, overlaid on the hydrostatic
  density --- the transonic isothermal-wind and hot-Parker families. A plausible
  wind therefore exists from step~0, which is needed where a cold start cannot
  launch (deep-RLOF or low-gravity boil-off). It is still EXHALE's own base with
  a thermodynamic overlay, so it stays consistent with the lower boundary; it
  only auto-enables \texttt{force\_start} so the heating can settle.
\item[\textbf{Starter} (starter solution / seed).] A \emph{complete, externally
  pre-converged} steady wind solution read in wholesale as the IC --- the
  Wind-AE warm start (\texttt{IC mode: windae}), which ramps a Wind-AE
  \texttt{windsoln} (or a grid/database seed) to this planet. Unlike a warm
  seed, a starter \emph{replaces the entire state} with another code's solution,
  including its own base; that base can clash with EXHALE's lower BC and pull the
  run off the wind branch (\S\ref{sec:bc-matrix}). A starter bootstraps hard
  cases but is not guaranteed consistent with EXHALE's boundary.
\end{description}

When a run does not restart from saved profiles
(\S\ref{sec:loadic}), the initial state is built by \texttt{set\_IC}
in \texttt{src/modules/init/}. Five options are available: the three IC
families \texttt{set\_IC} actually constructs (cold hydrostatic, transonic
isothermal wind, hot-Parker warm seed), plus \texttt{IC mode: auto}, which
selects among them at run time, and \texttt{IC mode: windae}, whose profile is
built in \texttt{init.f90} from the Wind-AE solver. The choice is made by
optional \texttt{input.inp} keys, with the cold hydrostatic atmosphere as the
default:

\begin{center}
\small
\begin{tabular}{@{}lll@{}}
\toprule
Family & Key & Default \\
\midrule
Cold hydrostatic + velocity seed & (none) & yes \\
Transonic isothermal wind & \texttt{Transonic IC: True} & off \\
Hot-Parker warm seed & \texttt{Hot Parker IC: $T_{\rm wind}$ [K]} & off \\
Automatic selection & \texttt{IC mode: auto} & off \\
Wind-AE warm start & \texttt{IC mode: windae} & off \\
\bottomrule
\end{tabular}
\end{center}

The transonic and hot-Parker families --- and only those two --- force the
first $1000$ iterations (\texttt{force\_start}) so that the heating structure
can develop before the convergence test is allowed to fire. Whatever the family,
the full dimensional IC is dumped to \texttt{output/IC\_dump.txt} for
inspection before the first step.

\paragraph{Automatic selection (\texttt{IC mode: auto}).}
\texttt{select\_IC\_auto} (\texttt{set\_IC.f90}) picks the family from a
single exact probe --- \texttt{find\_sonic} at the \emph{cold} base
sound speed in the actual potential. An interior sonic point selects
the transonic IC; this catches both deep-RLOF bases (the Roche
$\phi'\to0$ topology toward L1 always crosses the critical condition;
e.g.\ WASP-121b Case D, $r_c=1.30$) and low-gravity boil-off planets
($r_c \simeq b_0/2c_0^2$ falls inside the domain for small $b_0$).
Otherwise the cold hydrostatic default is kept (classic EUV-heated
winds: HD\,209458\,b has $r_c\sim54\,R_p$, far outside the domain).
There is no tunable threshold; $b_0$ (= the escape parameter $\Lambda$)
is logged as a regime diagnostic, the warm seed is never auto-selected
(the benchmark showed it buys nothing when the cold start works), and
the explicit keys override \texttt{auto}. Design rationale and
validation: \texttt{docs/auto\_ic\_design.md}.

\paragraph{Wind-AE warm start (\texttt{IC mode: windae}).}
With this key, \texttt{init.f90} builds the initial condition by running
the in-tree Fortran port of the Murray-Clay/Broome \emph{Wind-AE}
steady-state wind solver (\texttt{src/modules/wind\_ae/}). The bridge
\texttt{wae\_exhale\_bridge} maps the planet's EXHALE parameters
($M_{\rm p}$, $R_0$, $M_\star$, $a$, $F_{\rm XUV}=J_{\rm XUV}$, $L_\star$,
He/H, domain mode) to Wind-AE, ramps a shipped seed solution
(\texttt{inputdata/windae\_seed.csv}) to the planet by continuation, runs
the BVP relaxation $+$ outward integration, and writes
\texttt{output/\{Hydro\_ioniz,Ion\_species\}\_IC.txt} on the EXHALE grid;
\texttt{load\_IC} then ingests them. The whole EXHALE-input
$\to$ Wind-AE solve $\to$ IC $\to$ run is a single \texttt{EXHALE.x}
invocation. This is a metal-free H/He \emph{warm start} (EXHALE
re-equilibrates the ionization and any metals from the first step);
finish with \texttt{Solver: Newton} for a quantitative $\dot M$. Best
suited to seed-adjacent hot Jupiters; far-from-seed planets may fail the
ramp (the run then advises \texttt{IC mode: auto}). A standalone
generator \texttt{wind\_ae\_ic.x} (\texttt{make wind\_ae\_ic}) does the
same out of process. Full port/validation:
\texttt{md\_backup/wind\_ae\_fortran\_port\_plan.md}.

The guiding principle in all three cases is \emph{consistency with the
lower boundary condition} of \S\ref{sec:lower}: every IC matches
$(\rho, p, v) = (\rho_{\rm bc},\, n_{\rm tot,bc}{+}\delta p_{\rm bc},\, 0)$ at the
base exactly, so the boundary never has to fight the initial state.
The quantitative comparison of the families is in
\texttt{docs/initial\_condition\_benchmark.pdf}; the conclusions
quoted below come from there.

\subsection{Default: cold neutral hydrostatic atmosphere + velocity seed}
\label{sec:ic-hydrostatic}

The original ATES initial condition, kept as the EXHALE default
because it is the most robust (\S\ref{sec:ic-hotparker}).

\paragraph{Density.}
An isothermal hydrostatic atmosphere in the (Roche or spherical)
potential, anchored to the base density BC:
\begin{equation}
\rho(r) = \rho_{\rm bc}\,
   \exp\!\Big[\, b_{0,\rm eff}\big(\phi(r_0) - \phi(r)\big) \Big],
\label{eq:ic-hydro}
\end{equation}
with one deliberate distortion: $b_{0,\rm eff}$ starts at the physical
value ($1$ in code units, the potential already carrying the Jeans
parameter $b_0$) and is increased in steps of $0.2$ until the density
at the domain midpoint $r_{1/2} = (1 + r_{\max})/2$ falls below
$10^{-8}$ (\texttt{set\_IC.f90}, lines 85--101). For low-gravity
planets the true isothermal atmosphere would not decay enough across
the domain; the loop artificially \emph{steepens} the initial
atmosphere so that the outer half of the grid starts effectively
empty. The profile is then floored at $\rho = 10^{-8}$ (line 112).
Equation~(\ref{eq:ic-hydro}) with $b_{0,\rm eff} > 1$ is therefore
\emph{not} an equilibrium --- it is a concentrated cold reservoir that
the subsequent heating will evaporate.

\paragraph{Velocity.}
A small outward linear seed,
\begin{equation}
v(r) = \tfrac12\,(r - r_0),
\end{equation}
zero at the base and growing outward (line 108). Starting from exact
rest delays the launch of the wind; the seed gives the outflow a
direction without imposing any particular magnitude.

\paragraph{Pressure, temperature, ionization.}
Isothermal at the base temperature and consistent with the lower BC
(lines 137--146):
\begin{equation}
p = (n_{\rm tot,bc} + \delta p_{\rm bc})\,\frac{\rho}{\rho_{\rm bc}}, \qquad
T = 1 \;(= T_0\ \mbox{everywhere}),
\end{equation}
and a mostly neutral composition: $x_{\rm HII} = \delta p_{\rm bc}
\sim 10^{-10}$, $x_{\rm HeII} = 10^{-10}\,Y$, HeIII and He~$2^3$S
zero, the rest neutral. At the base cell this reproduces
$(\rho_{\rm bc},\, n_{\rm tot,bc}{+}\delta p_{\rm bc},\, T_0)$ exactly.

The physical picture: \emph{a cold, neutral, hydrostatically stratified
atmosphere is exposed to the stellar EUV flux at $t=0$ and left to
evaporate}. The first ionization-equilibrium sweep deposits the
heating, the upper layers warm and expand, and the wind launches
self-consistently --- typically within the forced-start phase.

\subsection{Transonic isothermal-wind IC (EXHALE)}
\label{sec:ic-transonic}

Motivation: deep Roche-lobe-overflow geometries (the canonical case is
WASP-121b Case~D, whose sonic point sits essentially at L1) launch
\emph{no} wind from the hydrostatic IC --- the run stalls in a bound,
breathing state. For these cases the run must start from a flowing
solution. With \texttt{Transonic IC: True}, \texttt{wind\_profile}
(\texttt{set\_IC.f90}, lines 177--233) solves the steady isothermal
wind in the actual ATES potential via its Bernoulli integral:
\begin{equation}
\tfrac12 v^2 - c^2 \ln v - 2c^2 \ln r + \phi(r) = B,
\qquad
c^2 = \frac{n_{\rm tot,bc} + \delta p_{\rm bc}}{\rho_{\rm bc}},
\label{eq:bernoulli}
\end{equation}
where $c$ is the dimensionless isothermal sound speed of the cold base
gas. The constant $B$ is fixed at the sonic point $r_c$, located by
bisection on the critical condition
\begin{equation}
\phi'(r_c) = \frac{2c^2}{r_c}
\end{equation}
(\texttt{find\_sonic}; in Roche mode $\phi' \to 0$ toward L1, so a
sign change is generic). At each radius Eq.~(\ref{eq:bernoulli}) has
one subsonic and one supersonic root; the transonic wind takes the
subsonic branch for $r < r_c$ and the supersonic branch for $r > r_c$
(robust bisection per cell, \texttt{wind\_root}). The density then
follows from steady mass conservation pinned to the base BC,
\begin{equation}
\rho v r^2 = \rho_{\rm bc}\, v(r_0)\, r_0^2 .
\end{equation}

If no interior sonic point exists (gravity dominates out to
$r_{\max}$: a bound ``breeze'' case), the routine reports failure and
\texttt{set\_IC} falls back to the hydrostatic IC of
\S\ref{sec:ic-hydrostatic} with a console notice. If the base itself
is already super-critical, the flow is taken supersonic from the base.

\subsection{Hot-Parker warm-seed IC (EXHALE)}
\label{sec:ic-hotparker}

This family is the product of the initial-condition benchmark
(\texttt{docs/initial\_condition\_benchmark.pdf}, \S7), whose central
negative result deserves restating: \textbf{seeding the full
isothermal Parker wind --- its density \emph{and} velocity --- makes
the run diverge.} The hot, low-density Parker column is globally
inconsistent with the non-isothermal ATES equilibrium; it drains onto
the base and drives the self-\emph{amplifying} breathing mode of
\S\ref{sec:breathing} (whereas the cold hydrostatic IC breathes
self-limitingly and converges). The lesson: the base may not be handed
an initial state whose global column it cannot support.

The robust compromise, enabled by
\texttt{Hot Parker IC: $T_{\rm wind}$} (default
$T_{\rm wind} = 10^4$~K), is a \emph{warm seed}: keep the stable cold
hydrostatic \emph{density} of \S\ref{sec:ic-hydrostatic} and overlay
only the fields whose relaxation is slow and stiff --- temperature,
ionization, and a velocity head-start. The overlay is anchored to the
cold base by a spatial smoothstep window (lines 114--133),
\begin{equation}
s(r) = \mathrm{clip}\!\left(\frac{r - 1}{r_{\rm tr} - 1},\, 0,\, 1\right),
\qquad
\xi(r) = s^2 (3 - 2s),
\end{equation}
with $r_{\rm tr} = $ \texttt{hp\_base\_rtr} $= 1.3\,R_p$ (a fixed default
in \texttt{parameters.f90}; it is a module variable, not a Fortran
\texttt{parameter}, but no input key exposes it), so that $\xi = 0$
at the base and $\xi = 1$ for $r \ge r_{\rm tr}$. The overlaid fields
are
\begin{align}
T(r)   &= 1 + \left(\tfrac{T_{\rm wind}}{T_0} - 1\right)\xi, &
v(r)   &= v_{\rm P}(r)\, s, \nonumber\\[2pt]
p(r)   &= \rho\,\tfrac{n_{\rm tot,bc}+\delta p_{\rm bc}}{\rho_{\rm bc}}\,
          T\,(1 + \xi), &
x_{\rm HII},\, x_{\rm HeII} &: \mbox{neutral} \to \mbox{ionized as } \xi,
\end{align}
where the $(1+\xi)$ factor in the pressure is the electron partial
pressure of the ionized warm gas (particle number doubles at full
ionization), and $v_{\rm P}$ is the Parker velocity head-start: the
transonic profile of \S\ref{sec:ic-transonic} re-solved with the
\emph{warm, ionized} sound speed
\begin{equation}
c^2_{\rm warm}
   = \frac{n_{\rm tot,bc}+\delta p_{\rm bc}}{\rho_{\rm bc}}\,
     \frac{T_{\rm wind}}{T_0}\times 2,
\end{equation}
its \emph{density discarded} (that is the part that destabilizes), and
its velocity multiplied by $s$ so it vanishes at the base. If the warm
profile has no interior sonic point, the head-start quietly degrades
to the small linear seed.

At the base $\xi = 0$ gives $T = T_0$, $p = n_{\rm tot,bc} + \delta p_{\rm bc}$,
$v = 0$ --- the lower BC of \S\ref{sec:lower} exactly, with no
inverted gradients for the first step to undo.

\paragraph{Benchmark verdict.}
The warm seed converges to the same $\dot M$ as the cold IC with no
NaN/blow-up, but yields \emph{no speedup} (${-}2.3\%$ steps in the
benchmark): ATES convergence is gated by the slow momentum ($du$)
relaxation, not by the thermal/ionization relaxation the warm seed
shortcuts. The cold hydrostatic IC therefore remains the default; the
warm seed is kept for cases where the cold start struggles to launch.

\subsection{Metals and common finalization}

In every family the metals start maximally neutral: the entire
elemental abundance is placed in the neutral ground stage,
\begin{equation}
f_{X^0} = \frac{n_X/n_{\rm H}}{1 + 4Y}, \qquad
f_{X^{+}} = f_{X^{++}} = 0
\end{equation}
(\texttt{set\_IC.f90}, lines 152--162; abundances from
\texttt{metals.inp}). The first ionization-equilibrium sweep
redistributes each element over its stages immediately, so the choice
costs at most one equilibrium solve. The complete IC (dimensional
$r, n, v, p, T$) is written to \texttt{output/IC\_dump.txt} before the
first step.

\subsection{Restarts (\texttt{Load IC})}
\label{sec:loadic}

When \texttt{Load IC} is enabled, \texttt{set\_IC} is bypassed and the
previous run's profiles (including ghost cells) are read back from the
\texttt{*\_IC.txt} copies (\texttt{load\_IC.f90}). This only sets the
initial state; \texttt{Apply\_BC} overwrites the ghost cells from the
restored interior at the first step, so a restart cannot smuggle in a
different boundary condition. (Operationally: the \texttt{output/}
$\to$ \texttt{*\_IC.txt} copy is performed by the GUI, not by the
Fortran code --- for scripted runs the \texttt{*\_IC.txt} files must
be seeded by hand, or the code stops on the empty files it creates.)

\paragraph{2026-08-13: a restart never carries zero metals into a
metals-on base.} The schema-2 writer emits the metal columns
unconditionally, so testing only for the \emph{presence} of a column
accepted a file written by a metals-off run and let the restart proceed
with zero metal density everywhere --- while the lower boundary still
counts the metals in the mass and particle budget under
\texttt{eos\_metals}, leaving an $O(1)$ residual in the first cells.
\texttt{load\_IC} now treats an element as absent in either case: no
column for some stage, \emph{or} columns that sum to zero over the
interior. Such an element is rebuilt exactly as a cold start builds it,
$n_X = a_X\,n_{\rm H}$ with all of it neutral, and the substitution is
reported; the rebuild happens \emph{before} \texttt{calc\_rho}
reconstructs the mass density, so the restart keeps the same mass closure
the cold start has by construction. The elements that were substituted are
recorded in \texttt{melem\_from\_abundance} and echoed by the setup report.

\paragraph{2026-08-15: a restart requires an IC written at the same $N$.}
Since the cell count is a runtime key (\S\ref{sec:overview}), an IC file
written at a different $N$ would previously have died in the middle of the
read with a bare end-of-file error. \texttt{load\_IC} now counts the data
records of \texttt{output/Hydro\_ioniz\_IC.txt} first --- the files carry
one record per cell, $N + 2N_g$ rows --- and stops with an explicit message
naming both counts when they disagree, so a mismatched
\texttt{Grid cells:} and a truncated file are reported rather than
guessed at. There is no interpolation onto a new grid: a restart requires
an IC written at the same $N$.

\section{Summary}

\begin{table}[h]
\centering
\small
\begin{tabular}{@{}llll@{}}
\toprule
Quantity & Lower boundary (base, $r=1$) & Upper boundary ($r=r_{\max}$) \\
\midrule
Density $\rho$       & Dirichlet: $\rho_{\rm bc} = (m_{\rm gas}/{\rm H})/(1+Y)$ (nuclei $=n_0$) & copy / linear extrapolation \\
Velocity $v$         & one-way valve $\max(v_1,0)$ (smooth / mass-flux opt.) & copy / linear extrapolation \\
Pressure $p$         & $n_{\rm tot,bc}+\delta p_{\rm bc}$, $\delta p_{\rm bc} = n_e(\mbox{base})/n_0$ & copy / linear extrapolation \\
Temperature $T$      & implied: $T = T_0$ exactly (\S\ref{sec:dpbc}) & implied by extrapolated $(\rho,p)$ \\
Ionization, metals   & equilibrium solve in ghosts (no BC) & equilibrium solve in ghosts (no BC) \\
Radiation            & in radiative shadow (large column) & unattenuated stellar spectrum \\
Type                 & isothermal--isochoric reservoir + valve & supersonic free outflow \\
Code                 & \texttt{Apply\_BC.f90:31--80} & \texttt{Apply\_BC.f90:95--99} \\
\bottomrule
\end{tabular}
\caption{Boundary conditions at a glance, with every option off (the
defaults). Ghost cells: $N_g = 2$ per
side; BCs act on ghosts only and reach the interior through the HLLC
fluxes. $m_{\rm gas}/{\rm H} = 1 + 4Y + \sum_X a_X A_X$ is the gas mass
per H nucleus from \texttt{composition.f90}, Eq.~(\ref{eq:rhobc}); it
reduces to the legacy $1+4Y$ for a metals-off atomic base. The velocity
and pressure rows have the opt-in alternatives of
\S\ref{sec:stabilizers}: mass-flux base velocity, and a ghost pressure
set by $\mathrm{d}p/\mathrm{d}r$ continuity (\texttt{Hydrostatic base})
or by $\mathrm{d}T/\mathrm{d}r = 0$ (\texttt{Base ghost temperature:
continuous}), the latter two replacing the implied $T = T_0$.}
\label{tab:summary}
\end{table}

The lower boundary is best read as a single physical statement ---
\emph{the wind is rooted in an infinite isothermal reservoir at
$(n_0, T_0)$ that supplies but never absorbs mass} --- implemented as
fixed $\rho$, electron-corrected fixed $p$, and a one-way velocity
valve. Its characteristic dynamical signature is the self-limiting
breathing oscillation of \S\ref{sec:breathing}, which is excluded
from convergence diagnostics via $r_{\rm esc}$, smoothed for the
Newton solver via \texttt{Valve eps}, and guarded against in the
advection post-processor. The upper boundary is a plain supersonic
free-outflow extrapolation whose validity rests on the wind crossing
the sonic point inside the domain.

The initial conditions are built to respect the same reservoir: all
three families (\S\ref{sec:ic}) match the base state
$(\rho_{\rm bc},\, n_{\rm tot,bc}{+}\delta p_{\rm bc},\, v{=}0)$ exactly, and the
benchmark experience condenses into one rule --- \emph{the density
column must be one the base can support}. The cold hydrostatic
atmosphere (default) satisfies this trivially and evaporates into the
wind on its own; the transonic IC replaces it with a flowing solution
where the hydrostatic start cannot launch (deep-RLOF, sonic point at
L1); the hot-Parker warm seed may overlay temperature, ionization, and
velocity on the cold density, but seeding the Parker \emph{density}
itself tips the breathing base from self-limiting to self-amplifying
and crashes the run.

\section{Empirical base-BC test matrix (HD209458b)}
\label{sec:bc-matrix}

\subsection{Convergence criterion: history and the flux-based decision}
\label{sec:conv-criterion}

The convergence test used here changed over the course of this work, and it is
worth stating plainly where it landed and why.

\begin{itemize}
\item \textbf{Before the Newton finish --- flux criterion.} The original
  time-marching run (PLM, then two-stage PLM$\to$WENO3) was stopped on the
  \emph{flux} criterion: the fractional variation of the mass flux,
  $\Delta\dot M/\dot M < 10^{-3}$, where $\dot M \propto \rho v r^2$ (the code's
  $du$ is precisely the radial spread of $\rho v r^2$ over the wind). This is the
  ATES criterion of Caldiroli et al.\ (2021) and is equivalent to the CETIMB
  requirement (Koskinen et al.\ 2013a) that the flux constant $F_c = \rho v r^2$
  be constant with altitude.
\item \textbf{With the Newton finish --- a stronger residual criterion.} When the
  JFNK steady solver was added, it was driven to a much stricter target: the
  steady residual $\lVert R\rVert = \lVert\mathrm{d}u/\mathrm{d}t\rVert < 10^{-3}$
  on \emph{every} conserved equation (mass, momentum, energy). This residual norm
  is far more demanding than flux flatness, and at the time these runs were made
  the JFNK solver reached it on only some of them; several full-physics runs were
  labeled ``JFNK fail'' even though their wind was perfectly flux-converged.
  \emph{(2026-08-10: the ``JFNK fail'' outcomes recorded in this section were
  produced with a diagonal scaling whose momentum floor was set by the base cell,
  and with a watchdog that aborted after 15 iterations without a new best
  $\lVert R\rVert$ --- before the opening pseudo-transient has finished raising
  it. Both have since been replaced and the same configurations converge; see
  \texttt{docs/newton\_scaling\_and\_base\_wall.md}. The residual-vs-flux
  \emph{decision} below is unaffected, but the ``held up by the near-base
  momentum'' attribution attached to individual rows is withdrawn.)}
\item \textbf{The reference codes use only the flux criterion.} Neither ATES
  (Caldiroli 2021: $\Delta\dot M/\dot M < 10^{-3}$) nor CETIMB (Koskinen 2013a:
  $F_c$ constant with altitude) imposes a residual-norm threshold. The reference
  standard for ``the wind has converged'' is flux flatness; the residual norm is
  an EXHALE addition tied to the optional Newton finish.
\item \textbf{Decision (adopted).} The \textbf{convergence is judged by the flux
  criterion}, matching the reference codes; the steady residual $\lVert R\rVert$
  is computed and reported \emph{for reference only} and gates the stop solely
  when the user explicitly requests it (\texttt{Resid tol:}). In the code the
  default monitor is flux-based ($du < du_{\rm th}$ plus mass-flux level
  stability) and the residual --- volume-weighted by default
  (\texttt{resid\_vol} in \texttt{parameters.f90}), since the L-$\infty$ max is
  inflated by the small near-base cells --- is printed each \texttt{N\_resid}
  steps as a diagnostic. The \texttt{verdict} column of the table below
  therefore reads \texttt{flux-conv}, not ``JFNK fail'', for runs that are
  flux-converged but whose residual finish did not reach the target.
\item \textbf{2026-08-11: the $du$ triggers are armed on an ascending crossing.}
  $du$ is the radial spread of $\rho v r^2$, so a \emph{small} $du$ is evidence
  of relaxation only for a state the marching loop has actually relaxed. A
  freshly generated IC --- cold hydrostatic, transonic, Wind-AE, hot Parker ---
  is an analytic profile whose $du$ measures the smoothness of the formula
  rather than the flatness of a wind: the WASP-121\,b transonic IC measured
  $du(1) = 1.85\times10^{-5}$, which tripped both the Newton hand-off and the
  secondary-ionization flip on step~1 and ran away to NaN. Each $du$-keyed
  trigger --- the $du < du_{\rm th}$ stop, and the JFNK hand-off that the
  secondary-ionization flip shares --- is therefore \emph{armed} only once $du$
  has been seen at or above its own threshold, and fires only on the way back
  down (\texttt{du\_stop\_armed}, \texttt{du\_newton\_armed} in
  \texttt{EXHALE\_main.f90}). A state read back with \texttt{Load IC? True} was
  relaxed by the marching loop that wrote it, so its $du$ is meaningful and both
  triggers start armed. Cold hydrostatic starts arm on step~1 anyway
  (measured $du(1) = 0.37$ for the WASP regression cases), so their behavior is
  unchanged. Details: \texttt{docs/newton\_scaling\_and\_base\_wall.md}, \S9.
\item \textbf{2026-08-15: the Wind-AE IC measured under that guard.} The
  arming guard had been measured on the transonic and cold-hydrostatic ICs but
  not on \texttt{IC mode: windae}, whose ${\sim}3$-step false $du$ stop is what
  originally forced the ``always Newton-finish'' rule. Measured on
  \texttt{examples/12\_windae\_ic\_hd209} with a 2000-step cap:
  $du$(step 2) $= 1.43\times10^{-4}$, below every threshold on the untouched
  generated IC, and nothing fires; the $du$ stop arms at step~23
  ($du = 1.027\times10^{-3}$) and the Newton hand-off at step~222
  ($du = 1.002\times10^{-2}$), both on the first ascending crossing. The false
  stop is structurally gone for this IC mode too. This does \emph{not} retire
  the rule of thumb: a $du$-stop $\dot M$ still carries the usual
  path-dependent spread, so the Newton finish stays the prescription for a
  quantitative rate.
\end{itemize}

\paragraph{2026-08-11: the tabulated rates are a product of the solver and the
cooling as they stood when the matrix was run.} Two changes since then move the
numbers in the two tables below, and neither has been re-measured on this matrix.
(i) The JFNK line search now accepts, stores, and judges convergence on the true
residual --- with the WENO smoothness weights recomputed at each trial state
rather than frozen at the previous iterate --- so runs that were recorded here as
failing their residual finish are not evidence about the current solver. (ii) The
two ground-term fine-structure coolants of the time --- the set has since grown to
the eight C\,\textsc{i}/C\,\textsc{ii}/N\,\textsc{ii}/O\,\textsc{i} lines --- are
allowed to be optically thick, and
the CHIANTI coronal fits are switched off smoothly below their $10^3$\,K validity
floor; on a planet whose base falls out of the coronal regime this changes the base
equilibrium temperature by hundreds of kelvin. The tables are kept as the dated
record of the base-BC comparison they were made for, and the \emph{relative}
ordering across base modes is what they are being cited for. They should not be
read as current mass-loss rates: the production HD~209458~b run under the present
code gives $\log_{10}\dot M = 9.47$ at $\lVert R\rVert = 1.8\times10^{-4}$ (the
2026-08-19 production set; it read $9.31$ at $7.7\times10^{-4}$ when this
paragraph was written on 2026-08-11). These
runs were not repeated, because repeating them is a re-measurement of the whole
matrix rather than a correction to it.

To pin down how the lower-boundary choices above actually affect a converged
wind, HD209458b was relaxed under a controlled combination of solver and base
options (the runs live in \texttt{EXHALE/HD209458b\_test/}). Each folder name
encodes the model as \texttt{<base>\_<solver>\_He\{0,1\}\_met\{0,CNO,FULL\}[\_flags]}:
\textbf{base} is \texttt{dens} ($n_0 = 10^{14}\,\mathrm{cm^{-3}}$ fixed; the
\texttt{dens12p5} run lowers it to $\log n_0 = 12.5$) or \texttt{pres1ub}
(pressure-anchored at $1\,\mu$bar, \S\ref{sec:stabilizers}); \textbf{solver} is
\texttt{PLM} (single-stage, stopped at a $du$ plateau), \texttt{PWN}
(two-stage PLM$\to$WENO3 then a Newton/JFNK finish), or \texttt{PLMN}
(single-stage PLM then JFNK, skipping WENO3); \textbf{S}/\textbf{V} are
the Shapiro filter / mass-flux base velocity of \S\ref{sec:stabilizers};
\textbf{He}/\textbf{met} are He~$2^3$S and the trace metals (C,N,O or the full
C,N,O,Mg,Ca,Na,Fe set); \textbf{IC} is the cold hydrostatic or the Wind-AE warm
start (\S\ref{sec:ic}). The results are split into a pressure-anchored-base
table (Table~\ref{tab:bc-pres}) and a density-anchored-base table
(Table~\ref{tab:bc-dens}).

In both tables: $\dot M$ is in g\,s$^{-1}$, $v_{\rm out}$ in km\,s$^{-1}$, and
neg.\% is the \emph{radius-weighted} infall fraction --- the fraction of the
linear radius range $[r_1, r_N]$ over which $v<0$ (each cell weighted by its
width $\Delta r$). This is used in place of the raw cell-count fraction, which
the non-uniform grid (cells packed near the base) inflates: the dense near-base
sliver of sub-$10\,\mathrm{m\,s^{-1}}$ noise would otherwise read as tens of
percent. The radius-weighted measure cleanly separates that base noise
($\lesssim 4\%$) from genuine infall (tens of percent). ``---'' marks a run on
an infall
branch (wind $\dot M$ undefined).

The \textbf{verdict} column reports convergence \emph{by the flux criterion} (the
$\Delta\dot M$\% column), \emph{not} by the residual: \texttt{Newton} $=$ a
genuine Newton-converged steady state (solver \texttt{info}$=$0,
$\lVert R\rVert<10^{-3}$); \texttt{flux-conv} $=$ flux-converged
($\Delta\dot M$\%$\,<1$) but the JFNK residual finish did not reach
$\lVert R\rVert<10^{-3}$ --- by the ATES/CETIMB standard these \emph{are} converged
winds, not failures; \texttt{marginal} $=$ $\Delta\dot M$\%$\,\sim$1--5\% (a wind,
not tightly flat; carries the $\sim$2--5\% path-dependent $\dot M$ spread of the
PLM-stall runs); \texttt{not conv} $=$ $\Delta\dot M$\%$\,>5\%$ (an unconverged /
breathing wind); \texttt{infall} $=$ relaxed to an accretion branch.

$\Delta\dot M$\% is the radial spread of the mass flux $\rho v r^2$ over the wind
region ($r \ge r_{\rm esc}$), $(\max-\min)/\mathrm{mean}$ --- this is the actual
\textbf{convergence criterion of the reference codes}: ATES (Caldiroli 2021)
stops at $\Delta\dot M/\dot M < 10^{-3}$, and CETIMB (Koskinen 2013a) requires
$\rho v r^2 = F_c$ to be constant with altitude. We mark $\lesssim1\%$ as
\emph{flux-converged} (a genuine steady wind \emph{by the reference standard},
independent of the $\lVert R\rVert$/Newton outcome); ``---'' = infall (not a
wind). Note that all \texttt{PWN} rows are flux-converged ($\le 0.9\%$) even
where the stricter $\lVert R\rVert$ test labels them ``JFNK fail''. The
``JFNK fail'' labels date from the pre-2026-08-10 solver; the flux verdicts do
not depend on them.

\begin{landscape}
\begin{table}[p]
\centering
\footnotesize
\setlength{\tabcolsep}{4pt}
\renewcommand{\arraystretch}{1.2}
\resizebox{\linewidth}{!}{%
\begin{tabular}{l l l c c c c l r r r r l}
\toprule
folder & BC & solver & S & V & He & met & IC & $\log\dot M$ & $v_{\rm out}$ & neg.\% & $\Delta\dot M$\% & verdict \\
\midrule
\texttt{pres1ub\_PWN\_He0\_met0\_cold}        & pres & PWN & \N & \N & \N & \N & cold & 10.39 & $+8.59$ & 0.0 & 0.53 & Newton \\
\texttt{pres1ub\_PWN\_He1\_metCNO\_cold}      & pres & PWN & \N & \N & \Y & CNO & cold & 9.53 & $+6.63$ & 1.6 & 0.62 & flux-conv \\
\texttt{pres1ub\_PWN\_He1\_metFULL\_cold}     & pres & PWN & \N & \N & \Y & FULL & cold & 9.37 & $+6.29$ & 2.2 & 0.87 & flux-conv \\
\texttt{pres1ub\_PWN\_He1\_metFULL\_cold\_duJFNK} & pres & PWN$^{(d)}$ & \N & \N & \Y & FULL & cold & 9.37 & $+6.29$ & 1.7 & 0.88 & flux-conv$^{(d)}$ \\
\texttt{pres1ub\_PLMN\_He1\_metFULL\_cold}    & pres & PLMN & \N & \N & \Y & FULL & cold & 9.40 & $+6.10$ & 1.8 & 3.5 & marginal$^\S$ \\
\texttt{pres1ub\_PWN\_He1\_metCNO\_cold\_Vflux} & pres & PWN & \N & \Y & \Y & CNO & cold & 9.53 & $+6.63$ & 1.6 & 0.74 & flux-conv \\
\texttt{pres1ub\_PWN\_He1\_metCNO\_cold\_Shap} & pres & PWN & \Y & \N & \Y & CNO & cold & 10.21 & $+7.45$ & 25.3 & 114 & no conv.$^\dagger$ \\
\texttt{pres1ub\_PLM\_He0\_met0\_cold}        & pres & PLM & \N & \N & \N & \N & cold & 10.40 & $+8.36$ & 0.0 & 3.5 & marginal \\
\bottomrule
\end{tabular}%
}
\caption{\textbf{Pressure-anchored base} ($1\,\mu$bar, $n_0\approx5\times10^{12}\,
\mathrm{cm^{-3}}$). Columns and general notation as described in the text.
$^\dagger$The pressure $+$ Shapiro run never converged: it ran to the
$10^6$-step cap with $25\%$ of the radius in infall (vs.\ $1.6\%$ for the same
run without Shapiro), so the filter both prevents convergence and deepens the
infall; the quoted $\dot M$ is the outer-region value, not a steady state.
$^\S$\texttt{PLMN} $=$ single-stage PLM $\to$ JFNK (the WENO3 stage skipped). It
reaches an \emph{essentially identical final state} to the \texttt{PWN}
(PLM$\to$WENO3$\to$JFNK) row above --- $\log\dot M = 9.40$ vs.\ $9.37$,
$v_{\rm out} = +6.10$ vs.\ $+6.29$, overlapping $T/v/n/p$ profiles --- in about
\emph{half} the steps ($9.4\times10^4$ vs.\ $1.7\times10^5$); both still fail
JFNK and stall, so the WENO3 stage is not the JFNK bottleneck.
$^{(d)}$The \texttt{\_duJFNK} row is the same full-metal case but with the
\emph{flux}-keyed JFNK hand-off (\texttt{Solver: Newton}, JFNK fired once the
wind reached $du<10^{-2}$ rather than $\lVert R\rVert<0.05$). With this hand-off
\emph{plus} the volume-weighted residual norm the JFNK drove $\lVert R\rVert$ down
to $2.3\times10^{-3}$ --- a $\sim$100$\times$ improvement over the $0.17$ it
reached with the $\lVert R\rVert$-keyed hand-off --- but it stalled there
($\lVert R\rVert$ would not drop below $\approx2\times10^{-3}$ for 15 iterations,
the holdout being the base cell $j{=}1$ momentum), so it still falls back and the
\emph{final wind is identical} to the \texttt{metFULL} PWN row above ($\dot M$,
$v_{\rm out}$, $\Delta\dot M$\% all unchanged). It confirms that the only barrier
to a genuine $du<10^{-3}$ / $\lVert R\rVert<10^{-3}$ convergence for full physics
is the localized near-base momentum (\S\ref{sec:breathing}), not the solver path.}
\label{tab:bc-pres}
\end{table}
\end{landscape}

\begin{landscape}
\begin{table}[p]
\centering
\footnotesize
\setlength{\tabcolsep}{4pt}
\renewcommand{\arraystretch}{1.2}
\resizebox{\linewidth}{!}{%
\begin{tabular}{l l l c c c c l r r r r l}
\toprule
folder & BC & solver & S & V & He & met & IC & $\log\dot M$ & $v_{\rm out}$ & neg.\% & $\Delta\dot M$\% & verdict \\
\midrule
\texttt{dens\_PWN\_He0\_met0\_cold}           & dens & PWN & \N & \N & \N & \N & cold & 10.50 & $+8.80$ & 0.0 & 0.41 & Newton \\
\texttt{dens12p5\_PWN\_He0\_met0\_cold}       & dens & PWN & \N & \N & \N & \N & cold & 10.38 & $+8.58$ & 0.0 & 0.42 & Newton \\
\texttt{dens\_PWN\_He1\_metCNO\_cold}         & dens & PWN & \N & \N & \Y & CNO & cold & 9.61 & $+6.79$ & 2.3 & 0.54 & flux-conv \\
\texttt{dens\_PWN\_He0\_met0\_windae\_spherical} & dens & PWN & \N & \N & \N & \N & windae & 10.37 & $+5.77$ & 1.4 & 55 & not conv$^\ddagger$ \\
\texttt{dens\_PLM\_He0\_metCNO\_cold}         & dens & PLM & \N & \N & \N & CNO & cold & 9.72 & $+6.77$ & 3.9 & 4.2 & marginal \\
\texttt{dens\_PLM\_He0\_met0\_cold\_Vflux}    & dens & PLM & \N & \Y & \N & \N & cold & 10.50 & $+8.58$ & 0.0 & 3.5 & marginal \\
\texttt{dens\_PLM\_He0\_met0\_cold\_Shap}     & dens & PLM & \Y & \N & \N & \N & cold & --- & $-13.34$ & 100 & --- & infall \\
\texttt{dens\_PLM\_He0\_met0\_cold\_Shap\_Vflux} & dens & PLM & \Y & \Y & \N & \N & cold & --- & $-13.32$ & 100 & --- & infall \\
\texttt{dens\_PLM\_He0\_met0\_windae\_spherical} & dens & PLM & \N & \N & \N & \N & windae & 10.08 & $+3.68$ & 0.6 & 5.1 & not conv \\
\texttt{dens\_PLM\_He0\_met0\_windae\_spherical\_Shap\_Vflux\_Force} & dens & PLM & \Y & \Y & \N & \N & windae$+$F & --- & $-20.03$ & 95 & --- & infall \\
\bottomrule
\end{tabular}%
}
\caption{\textbf{Density-anchored base} ($n_0 = 10^{14}\,\mathrm{cm^{-3}}$;
\texttt{dens12p5}: $\log n_0 = 12.5$). Columns and general notation as described
in the text.
$^\ddagger$Wind-AE warm start: JFNK never triggered ($\lVert R\rVert$ stayed
above the switch); the Wind-AE IC marched under WENO3 drifts to an unconverged
breathing state ($du\!\sim\!0.7$), mostly outflow --- it does \emph{not} give a
clean Newton solution.}
\label{tab:bc-dens}
\end{table}
\end{landscape}

\paragraph{What the matrix shows.}
\begin{itemize}
\item \textbf{The Shapiro filter drives infall.} On the cold-IC density runs,
  every S-on case relaxes to a global infall ($v<0$ over $\approx100\%$ of the
  radius) while
  the matched S-off case is a clean outflow, independent of the velocity BC ---
  the headline reason the filter is opt-in (\S\ref{sec:stabilizers}). On the
  pressure $+$ full-physics case the filter instead prevents convergence
  altogether (row $\dagger$).
\item \textbf{The base BC barely changes the converged wind --- the lever is the
  base \emph{density}.} With a Newton finish, density ($n_0=10^{14}$) and
  pressure ($1\,\mu$bar) bases give nearly the same $\dot M$ (10.50 vs.\ 10.39).
  A \emph{density}-anchored run at the same low density (\texttt{dens12p5},
  $\log n_0 = 12.5$) reproduces the pressure result almost exactly (10.38 vs.\
  10.39; $+8.58$ vs.\ $+8.59$), confirming that pressure-anchoring acts only
  through the lower base density, not through the anchoring variable itself.
  The mass-flux velocity (V) is benign (it reproduces the valve result).
\item \textbf{Physics lowers $\dot M$.} He~$2^3$S $+$ C/N/O cooling drops
  $\log\dot M$ from $\sim$10.4 to $\sim$9.5--9.6; the full metal set
  (adding Mg/Ca/Na/Fe line cooling) lowers it further to $9.37$. With full
  physics the JFNK solve fails for both bases and falls back to a stall.
\item \textbf{Skipping WENO3 before the Newton finish (\texttt{PLMN}) changes
  nothing but the cost.} On the full-metal case, \texttt{PLMN} (PLM$\to$JFNK)
  reaches an essentially identical final state to \texttt{PWN}
  (PLM$\to$WENO3$\to$JFNK) --- $\log\dot M = 9.40$ vs.\ $9.37$, overlapping
  profiles --- in about half the steps, and both still fail JFNK and stall. So
  the reconstruction path is not the JFNK bottleneck; the stiff coupled
  (He~$2^3$S $+$ 7-metal) residual is. A genuine Newton finish for full physics
  would need a smoother residual (tighter chemistry tolerance) or a relaxed
  target, not a different reconstruction.
\item \textbf{The Wind-AE warm start does not help here.} It either du-stops at
  step~2 (returning $\approx$ the IC) or, with a Newton finish, never triggers
  JFNK and drifts to an unconverged breathing state. The cold hydrostatic IC,
  consistent with EXHALE's own base, is the reliable choice.
\end{itemize}

\subsection{Assessment: does the Wind-AE starter help?}
\label{sec:windae-assess}

\paragraph{So far, for HD209458b, the Wind-AE starter does not help --- the cold
IC is better.} Across every Wind-AE-seeded run in Table~\ref{tab:bc-dens} the
outcome is one of: a $du$-stop at step~2 that simply returns the seed (no real
relaxation), an unconverged breathing state when a Newton finish is requested
(JFNK never triggers), or, when marching is forced, a drift to infall
($v_{\rm out} = -20\,\mathrm{km\,s^{-1}}$). An earlier scratch run converged
instead to a global \emph{infall root}. None reached a clean, Newton-converged
outflow, whereas the cold IC with a \texttt{PWN} finish does (solver
\texttt{info}=0, clean wind).

\paragraph{Why.} A \emph{starter} replaces the entire state, \emph{including its
own base}, with another code's solution (\S\ref{sec:ic}). That imported base is
generally inconsistent with EXHALE's fixed lower boundary ($\rho_{\rm bc}$,
$T_0$); the mismatch perturbs the near-base flow off the transonic-wind saddle,
and since that saddle is unstable, the run is pulled to the nearby attractors ---
infall, or the breathing limit cycle. This is the same dense-base/lower-BC
tension that drives the breathing (\S\ref{sec:breathing}), now triggered by the
seed rather than by a bad global IC.

\paragraph{In fairness.} This does \emph{not} make Wind-AE useless. It was
brought in for the \emph{hard} cases where a cold start fundamentally cannot
launch a wind --- deep Roche-lobe-overflow geometries (e.g.\ WASP-121b Case~D,
with the sonic point essentially at L1) and low-gravity boil-off planets ---
where it is, at present, the only way to obtain a transonic seed. HD209458b is
an easy case for which the cold IC already works, so there the starter is both
unnecessary \emph{and} mildly harmful through the base mismatch. The honest
statement is therefore: \emph{the starter is unnecessary for easy cases and,
for its intended hard cases, has not yet delivered a clean solution because the
base-BC reconciliation is unsolved} --- not that the approach is without value.

\paragraph{What would make it useful.} The base mismatch must be removed. Two
plausible routes (untested): (i)~have the Wind-AE bridge hand off only a
\emph{warm overlay} --- temperature, ionization, and velocity --- on top of
EXHALE's own hydrostatic density and base, i.e.\ demote the ``starter'' to a
``warm seed'' (\S\ref{sec:ic}) that stays consistent with the lower boundary;
or (ii)~ramp the imported base onto EXHALE's $(\rho_{\rm bc}, T_0)$ boundary
before releasing the hydrodynamics. Either could keep the transonic-seed
advantage Wind-AE offers for the RLOF cases while avoiding the infall/breathing
seen here.

% ===================================================================== %
\section*{References}
\addcontentsline{toc}{section}{References}
% ===================================================================== %
\begingroup
\setlength{\parindent}{0pt}%
\setlength{\parskip}{4pt}%
\everypar{\hangindent=1.5em\hangafter=1}

Caldiroli, A., Haardt, F., Gallo, E., Spinelli, R., Malsky, I., \& Rauscher, E.
2021, A\&A, 655, A30. ``Irradiation-driven escape of primordial planetary
atmospheres.~I. The ATES photoionization hydrodynamics code.'' (The original
ATES code EXHALE is built on.)

Huang, C., Koskinen, T., Lavvas, P., \& Fossati, L. 2023, ApJ, 951, 123. ``A
Hydrodynamic Study of the Escape of Metal Species and Excited Hydrogen from the
Atmosphere of the Hot Jupiter WASP-121\,b.'' (The WASP-121\,b escape model used
here as the validation target; its description of CETIMB is what pointed to the
Shapiro-filter remedy ported in \S\ref{sec:stabilizers}.)

Koskinen, T.~T., Harris, M.~J., Yelle, R.~V., \& Lavvas, P. 2013a, Icarus, 226,
1678. ``The escape of heavy atoms from the ionosphere of HD\,209458\,b.~I. A
photochemical--dynamical model of the thermosphere.'' (CETIMB; the source of the
periodic Shapiro low-pass used to damp the gravity-unbalanced base sound waves.)

Shapiro, R. 1970, Reviews of Geophysics and Space Physics, 8, 359.
``Smoothing, filtering, and boundary effects.'' (The 1--2--1 low-pass filter
applied here.)
\endgroup

\end{document}
